\section{The Grammar Concept: Two-IR-Layer Model}
\label{sec:grammar-concept}

This section defines the ``grammar'' abstraction that structures the diagram compiler. A grammar is a self-contained visual vocabulary with its own Domain IR and layout engine(s), all compiling down to the shared Scene IR.

\subsection{What Is a Grammar?}

A \textbf{grammar} in this architecture is:

\begin{quote}
A grammar = (Domain IR, Layout Engine(s), Theme Extensions) that compiles a declarative description of \emph{what to communicate} into Scene IR primitives specifying \emph{what to draw}.
\end{quote}

Each grammar is:
\begin{itemize}
  \item \textbf{Self-contained}: It defines its own IR schema, validation rules, and layout logic.
  \item \textbf{Kernel-dependent}: It consumes kernel services (Scene IR, backends, themes, icons) but does not modify them.
  \item \textbf{Peer to other grammars}: Timeline, flow, and graph grammars are siblings, not nested.
\end{itemize}

% -------------------------------------------------------------------------
\subsection{Grounding in the Grammar of Graphics}
\label{sec:grammar-gog}
% -------------------------------------------------------------------------

\textbf{What the literature says.}
This architecture stands on the shoulders of three decades of
visualization-grammar research.

Wilkinson's \emph{Grammar of Graphics}~\cite{grammarofgraphics2005}
established that any statistical graphic can be decomposed into a small
algebra of orthogonal components: \textbf{data}, \textbf{aesthetic
mappings}, \textbf{geometric marks}, \textbf{statistical transforms},
\textbf{scales}, \textbf{coordinates}, and \textbf{facets}.
The key insight is generativity: a finite, composable set of primitives
describes an infinite variety of charts.
This is precisely the principle underlying the diagram compiler's grammar
abstraction --- a finite set of mark types (box, arrow, label, group,
connector) combines through composition to cover all target diagram classes.

Wickham's layered grammar~\cite{layeredgrammar2010} operationalized Wilkinson
into ggplot2, adding two engineering contributions of direct relevance:
(1)~\textbf{default inference} --- the system fills in unspecified dimensions
(axes, colors, font sizes) from data type and context, so the author
specifies only deviations from sensible defaults; and (2)~\textbf{spec/render
separation} --- a ggplot object is a spec; calling \texttt{ggsave()} triggers
the rendering pipeline.
Both principles are adopted verbatim: Domain IR fields have sensible defaults
inferred by the layout engine, and the render pass is strictly separated from
the compile pass.

Satyanarayan et al.'s Vega-Lite~\cite{vegalite2017} pushed the grammar further:
a concise, JSON-serializable spec compiles to Vega's lower-level dataflow IR,
which backends render to SVG or Canvas.
Vega-Lite's composition algebra (\texttt{layer}, \texttt{hconcat},
\texttt{vconcat}, \texttt{facet}) is the direct ancestor of the composition
layer in Section~\ref{sec:composition}.

\medskip\noindent
\textbf{The diagram extension.}
Where the GoG lineage addresses \emph{statistical graphics}, this architecture
addresses \emph{diagram types} --- swimlane timelines, flow pipelines, node-link
graphs, step-card sequences, and comparison matrices.
These diagrams have no \emph{data/aesthetic-mapping} duality; they have a
\emph{semantic entity / visual layout} duality instead.
A Timeline activity has \texttt{start}, \texttt{end}, and \texttt{track};
the compiler maps these to rectangle position, width, and vertical row ---
an analogous but distinct encoding to GoG's \texttt{field $\to$ channel}
mapping.

% -------------------------------------------------------------------------
\subsection{Agent-Authoring Reliability: Constrained DSL}
\label{sec:grammar-llm}
% -------------------------------------------------------------------------

\textbf{What the literature says.}
A growing body of constrained-generation research establishes that
\emph{grammar-constrained decoding is the essential reliability lever for
LLM-DSL pipelines}, and that small, well-constrained grammars lower the
failure surface dramatically.

Willard \& Louf~\cite{willard2023} show that reformulating LLM generation as
transitions in a finite-state machine derived from a formal grammar (or
regex) eliminates an entire class of syntax failures with negligible
per-token overhead.
Wang et al.~\cite{wang2023grammar} demonstrate that a \emph{minimal grammar
fragment} for a specific diagram instance is more reliable than the full
schema: ``grammar constraints derived from domain-specific knowledge enable
LLMs to generalize from few examples on complex DSL generation tasks.''
Dong et al.'s XGrammar~\cite{dong2024xgrammar} achieves up to $100\times$
speedup over naive per-step CFG execution by separating context-independent
from context-dependent tokens at grammar-compilation time, making
schema-constrained diagram generation practical in production LLM
serving stacks.

Tian et al.~\cite{tian2023chartgpt} studied LLM failure modes for chart spec
generation: LLMs produce syntactically valid but semantically wrong specs,
miss required fields, and use ambiguous field names.
Their mitigations map directly to the compiler's architecture: (1)~grammar
constraints eliminate syntactic failures; (2)~the validation layer (§\ref{sec:two-ir-model}) catches semantic failures; (3)~decomposed
generation (entity identification, then IR emission) handles complex diagrams.

Narechania et al.~\cite{narechania2021nl4dv} showed that targeting a
well-specified JSON grammar (Vega-Lite's schema) rather than a general-purpose
API reduces NL-to-visualization complexity substantially.
For small or on-device models, Ray~\cite{ray2026constraint} characterises the
``constraint tax'' --- hard schema-decoding improves syntactic validity but
can lower semantic accuracy for sub-3B models --- and recommends a
``reason free, constrain late'' strategy.

\medskip\noindent
\textbf{What we recommend.}
The Domain IR JSON Schema (one per grammar) is published as a formal
constraint and submitted to OpenAI Structured Outputs~\cite{openai-structured-outputs},
XGrammar~\cite{dong2024xgrammar}, or llama.cpp
GBNF~\cite{llama2024gbnf} for grammar-constrained generation.
Keeping each Domain IR \emph{small} (few required fields, flat nesting,
minimal enum values) is not a simplification concession --- it is a
reliability design decision grounded in the literature above.
The god-IR (Section~\ref{sec:two-ir-model}) is rejected in part because a
large, polymorphic schema is an agent-authoring failure vector.

\subsection{The Two-IR-Layer Model}
\label{sec:two-ir-model}

The architecture uses \textbf{two layers of intermediate representation}:

\ourdiagram[0.95\linewidth]{two-ir-model}{%
  The two-IR-layer model: each grammar's Domain IR feeds its own Layout Engine,
  all of which compile down to the single shared Scene IR,
  consumed by SVG, Skia, and PPTX rendering backends.}

\paragraph{Domain IRs are diverse, small, and grammar-specific.}
The Timeline IR has tracks, activities, milestones. The Flow IR has nodes, edges, lanes. The Graph IR has vertices, edges, clusters. These entities are \emph{semantically meaningful} within their grammar---an Activity knows it has a start date; a Node knows it has incoming edges.

Domain IRs are kept \textbf{small on purpose}:
\begin{itemize}
  \item A small grammar is reliably emittable by an LLM
  \item A small schema fits in a context window
  \item A small grammar has clear boundaries and fewer edge cases
\end{itemize}

\paragraph{Scene IR is the universal assembly language.}
All Domain IRs compile \emph{down} to the single shared Scene IR. The Scene IR knows nothing about timelines, flows, or graphs---it knows only primitives (Rect, Circle, Line, Path, Text, Group) and effects.

The Scene IR is the \textbf{one integration point}:
\begin{itemize}
  \item Backends consume only Scene IR
  \item Themes style Scene IR primitives
  \item Golden tests can verify Scene IR determinism
  \item Animation hints attach to Scene IR primitives
\end{itemize}

\subsection{Rejecting the God-IR Approach}

An alternative architecture would use a single ``god IR'' that tries to express all diagram types in one schema:

\begin{lstlisting}[language=yaml,caption={Hypothetical god-IR (rejected)}]
# DON'T DO THIS
diagram:
  type: timeline | flow | graph | comparison | ...
  entities:
    - kind: activity | node | cell | stat | ...
      properties: { ... massive union of all possible fields ... }
\end{lstlisting}

This approach is explicitly rejected because:

\begin{description}
  \item[Semantic muddiness] A single schema that means different things in different contexts is harder to validate, harder to understand, and harder to document.
  
  \item[Schema bloat] Every grammar's fields appear in the union, even when irrelevant. An activity has \texttt{start}/\texttt{end}; a graph node does not. The god-IR carries dead weight.
  
  \item[LLM hostility] A large, polymorphic schema with conditional validity rules is difficult for agents to emit correctly. Small, single-purpose schemas are easier.
  
  \item[Brittle evolution] Adding a new grammar requires modifying the central schema, risking regressions in all other grammars.
  
  \item[No clear ownership] Who owns the god-IR? Every grammar must negotiate changes. Separate Domain IRs have clear ownership.
\end{description}

The two-IR-layer model inverts the coupling: Domain IRs are independent; the Scene IR is the shared core. Adding a new grammar does not touch existing grammars' IRs.

\subsection{Composition IR: Above Domain IRs}

For multi-panel posters (Section~\ref{sec:composition}), a Composition IR sits \emph{above} the Domain IRs, not merged into them:

\ourdiagram[0.9\linewidth]{composition-ir}{%
  Composition IR sits above Domain IRs: the Composition Layout Engine dispatches
  each panel to its grammar-specific layout engine, then assembles the resulting
  sub-scenes into a single Unified Scene IR.}

The Composition IR \emph{references} panels, each with its own grammar type and Domain IR. It does not flatten them into a single schema.

\subsection{Grammar Interface Contract}

Every grammar must implement:

\begin{lstlisting}[language=javascript,caption={Grammar interface contract}]
interface Grammar<DomainIR> {
  // Schema
  readonly name: string;               // e.g., 'timeline', 'flow'
  readonly schema: JsonSchema;         // JSON Schema for DomainIR
  
  // Validation
  validate(ir: unknown): ValidationResult;
  
  // Layout
  layout(ir: DomainIR, theme: Theme): Scene;
  
  // Optional: grammar-specific theme extensions
  readonly themeSchema?: JsonSchema;
}
\end{lstlisting}

\paragraph{Schema.}
The JSON Schema for the Domain IR, used for validation and agent generation.

\paragraph{Validate.}
Checks the IR against schema and semantic rules. Returns diagnostics.

\paragraph{Layout.}
The core transformation: Domain IR + Theme → Scene IR. This is where all layout computation happens.

\paragraph{Theme extensions.}
Grammars may extend the theme schema with grammar-specific properties (e.g., \texttt{timeline.trackHeight}, \texttt{flow.nodeSpacing}).

\subsection{Current and Planned Grammars}

\begin{center}
\small
\begin{tabular}{llll}
\toprule
\textbf{Grammar} & \textbf{Domain Entities} & \textbf{Status} & \textbf{Family} \\
\midrule
Timeline & tracks, activities, milestones, sections & Implemented & Timeline/project \\
Flow & nodes, edges, lanes, legends & Implemented & Node-link/graph \\
Sequence & participants, messages, activations, fragments & Implemented & UML/software \\
Tree & root, children (recursive), labels & Implemented & Tree/hierarchy \\
\midrule
Class & classes, relationships, packages & Planned (Tier-1) & UML/software \\
State & states, transitions, composites & Planned (Tier-1) & UML/software \\
ER & entities, relationships, cardinalities & Planned (Tier-1) & UML/software \\
Chart & data, encoding, marks, scales & Planned (Tier-2) & Charts \\
\bottomrule
\end{tabular}
\end{center}

The full 22-type coverage matrix and family taxonomy are detailed in
Section~\ref{sec:family-taxonomy}.

\subsection{The Compounding Payoff}

The kernel/grammar separation creates a compounding payoff:

\begin{itemize}
  \item \textbf{Animation} added to the kernel is inherited by all grammars
  \item \textbf{New backends} (e.g., Keynote) serve all grammars immediately
  \item \textbf{Theme improvements} (new fonts, colors, effects) apply everywhere
  \item \textbf{Determinism verification} at the Scene level covers all grammars
  \item \textbf{Icon registry} additions are available to all grammars
\end{itemize}

Each new grammar starts with a full feature set: determinism, theming, SVG/PNG/PDF/PPTX output, animation support, validation, and agent integration---all for free.
